%% ========================================================================
%% Brāhmasphuṭasiddhānta of Brahmagupta (628 CE) — CHUNK 5
%% Pages 561–700 of the PDF (printed pages 171–315)
%% Chapters XIV–XXII (portions)
%% ========================================================================
\documentclass[11pt,a4paper]{article}

%% -------- Core Packages --------
\usepackage{fontspec}
\usepackage{polyglossia}
\usepackage{amsmath,amssymb,amsthm}
\usepackage{geometry}
\usepackage{graphicx}
\usepackage{xcolor}
\usepackage{titlesec}
\usepackage{fancyhdr}
\usepackage{enumitem}
\usepackage{array,booktabs,longtable,multirow}
\usepackage{verse}
\usepackage{lettrine}
\usepackage{microtype}
\usepackage{tocloft}
\usepackage{indentfirst}
\usepackage{etoolbox}

%% -------- Page Geometry --------
\geometry{
  paperwidth=210mm,paperheight=297mm,
  left=22mm,right=22mm,top=28mm,bottom=28mm,
  headheight=14pt,footskip=30pt
}

%% -------- Language Setup --------
\setmainlanguage{english}
\setotherlanguage{sanskrit}

%% -------- Font Configuration --------
\setmainfont{Linux Libertine O}
\setsansfont{Linux Biolinum O}
\setmonofont{DejaVu Sans Mono}

%% Devanagari (Sanskrit/Hindi) Font
\newfontfamily\devanagarifont[Script=Devanagari,Scale=1.15]{Noto Serif Devanagari}
\newfontfamily\devanagarifontsf[Script=Devanagari,Scale=1.1]{Noto Sans Devanagari}

%% -------- Colors --------
\definecolor{titlecolor}{RGB}{45,55,72}
\definecolor{sectioncolor}{RGB}{55,65,81}
\definecolor{notecolor}{RGB}{120,53,15}
\definecolor{rulecolor}{RGB}{180,160,140}
\definecolor{versecond}{RGB}{90,50,30}

%% -------- Section Formatting --------
\titleformat{\section}
  {\normalfont\Large\bfseries\color{sectioncolor}}
  {\thesection}{1em}{}
\titleformat{\subsection}
  {\normalfont\large\bfseries\color{sectioncolor}}
  {\thesubsection}{1em}{}
\titleformat{\subsubsection}
  {\normalfont\normalsize\bfseries\color{sectioncolor}}
  {\thesubsubsection}{1em}{}

%% -------- Header/Footer --------
\pagestyle{fancy}
\fancyhf{}
\fancyhead[L]{\small\textit{Brāhmasphuṭasiddhānta — Chunk 5 (Chs. XIV–XXII)}}
\fancyhead[R]{\small\thepage}
\renewcommand{\headrulewidth}{0.4pt}
\renewcommand{\footrulewidth}{0pt}

%% -------- Custom Commands --------
\newcommand{\longline}{\noindent\rule{\textwidth}{0.4pt}}
\newcommand{\shortline}{\noindent\rule{0.5\textwidth}{0.4pt}\par}
\newcommand{\cnote}[1]{\textcolor{notecolor}{\textit{#1}}}
\newcommand{\dev}[1]{{\devanagarifont #1}}
\newcommand{\dsec}[1]{\section*{\dev{#1}}}
\newcommand{\dsubsec}[1]{\subsection*{\dev{#1}}}

%% Verse environment for Sanskrit shloka
\newenvironment{shloka}{\begin{verse}\devanagarifont\large}{\end{verse}}

%% Footnote-style variant readings
\newcommand{\varreading}[1]{\textcolor{notecolor}{\small [#1]}}

%% -------- Hyperref Setup --------

%% -------- TOC Depth --------
\setcounter{tocdepth}{3}
\setcounter{secnumdepth}{3}

%% =======================================================================

\newcommand{\vs}{\vspace{0.5em}}
\begin{document}

%% ========================================================================
%% TITLE PAGE
%% ========================================================================
\begin{titlepage}
\centering
\vspace*{2cm}

{\Huge\bfseries\color{titlecolor} \dev{ब्राह्मस्फुटसिद्धान्त}}\\[0.8cm]
{\LARGE\color{sectioncolor} Brāhmasphuṭasiddhānta}\\[0.3cm]
{\large of \textbf{Brahmagupta} (628 CE)}\\[1.5cm]

\shortline

{\Large\bfseries\color{sectioncolor} Chunk 5: Pages 561–700}\\[0.3cm]
{\large Printed pages 171–315 of Volume I}\\[0.5cm]
{\normalsize Covering Chapters XIV–XXII (partial)}\\[1.5cm]

\shortline

\vspace{1cm}

{\small\textit{Digital edition prepared from}}\\
{\small The Indian Institute of Astronomical \& Sanskrit Research edition}\\[0.3cm]
{\small With corrections and variant readings from multiple manuscripts}\\[1cm]

\vfill

{\footnotesize Typeset in XeLaTeX with Noto Serif Devanagari}\\
{\footnotesize Sanskrit text digitized by Takao Hayashi (GRETIL, 1993)}\\
{\footnotesize Based on S.~Dvivedin's edition (Benares, 1902)}

\end{titlepage}

%% ========================================================================
%% TABLE OF CONTENTS
%% ========================================================================
\tableofcontents
\newpage

%% ========================================================================
%% INTRODUCTION
%% ========================================================================
\section*{Introduction}
\addcontentsline{toc}{section}{Introduction}

This volume contains pages 561–700 of the PDF of the
\textit{Brāhmasphuṭasiddhānta} (BSS) of Brahmagupta, corresponding to
printed pages 171–315 of Volume~I. This span covers the latter portion
of the astronomical chapters and the core mathematical chapters of the
treatise.

\subsection*{Overview of Chapters}
\addcontentsline{toc}{subsection}{Overview of Chapters}

The \textit{Brāhmasphuṭasiddhānta} consists of 24 chapters. The
chapters covered in this chunk are:

\begin{enumerate}[label=\textbf{Ch. \Roman*:}, wide=0pt]
  \item \textbf{Sphoṭagatyuttarādhikāraḥ} (Ch.~14) — On true planetary
    motions and their corrections. Covers the computation of true
    longitude from mean longitude using epicyclic corrections.
  \item \textbf{Tripraśnottarādhikāraḥ} (Ch.~15) — The three questions:
    relating to time, place, and direction; including shadow and gnomon
    calculations.
  \item \textbf{Grahāṇottarādhikāraḥ} (Ch.~16) — Advanced eclipse
    computations, including graphical methods (\textit{parilekha}).
  \item \textbf{Śṛṅgonnatyuttarādhikāraḥ} (Ch.~17) — On lunar latitude
    and related calculations.
  \item \textbf{Kuṭṭakādhyāyaḥ} (Ch.~18) — The algebra chapter. The
    longest mathematical chapter, covering the kuṭṭaka (pulverizer)
    method for solving linear Diophantine equations, operations with
    positive/negative numbers and zero, saṅkramaṇa, karaṇī (surds),
    and extensive algebra including the \textit{varga-prakṛti} (Pell's
    equation).
  \item \textbf{Śaṅku-chāyā-vijñānam} (Ch.~19) — On gnomon (śaṅku) and
    shadow (chāyā) computations; practical geometry and instrument
    usage.
  \item \textbf{Chandaś-citty-uttarādhikāyaḥ} (Ch.~20) — On Sanskrit
    prosody (chandah-śāstra), including combinatorial calculations for
    metrical patterns.
  \item \textbf{Golādhyāyaḥ} (Ch.~21, partial) — The sphere/cosmology
    chapter. This chunk contains the \textit{jyā-prakaraṇa} (sine
    computation) section (verses 17–23).
  \item \textbf{Yantrādhyāyaḥ} (Ch.~22) — On astronomical instruments
    (\textit{yantra}), including the \textit{dhanur-yantra} (bow
    instrument), \textit{yasti-yantra} (staff instrument), water
    instruments, and observation techniques.
\end{enumerate}

\subsection*{Source and Notes}
\addcontentsline{toc}{subsection}{Source and Notes}

The Sanskrit text presented here is based on multiple sources:
\begin{itemize}
  \item The GRETIL digitization (Takao Hayashi, 1993) based on
    S.~Dvivedin's edition (Benares, 1902), used for Chapters 18–21.
  \item Manual transcription from the scanned pages of the Indian
    Institute edition for Chapters 14–17 and 22.
  \item Variant readings (\textit{pāṭha-bhedaḥ}) are recorded in the
    footnotes as found in the printed edition, with manuscript
    sigla: (क) for the Poona manuscript, (ग) for the Bombay manuscript,
    (च) for the Baroda manuscript, and (छ) for the Hoshiarpur manuscript.
\end{itemize}

\longline

\newpage

%% ========================================================================
%% CHAPTER XIV: SPHOTAGATYUTTARADHIKARAH
%% ========================================================================
\dsec{अथ स्फुटगत्युत्तराध्यायः — चतुर्दशः}
\addcontentsline{toc}{section}{Chapter XIV: Sphoṭagatyuttarādhikāraḥ}

\cnote{[Printed pages 171–181. This chapter continues from the previous
chunk. The verses below begin from the portion available in the scanned
pages of this chunk.]}

\begin{shloka}
युगभागः कोटिज्यां कोट्यां करोति बाहुज्याम् ।\\
कोटिं भुजेन बाहुं क्षेत्रच्चावा स्फुटगतिकः सः ॥ १ ॥
\end{shloka}

\begin{shloka}
परमण्डलकर्णविकः करोति कोटिछाया स्फुटं कर्तुम् ।\\
कर्णान्तकोटिकोण्डा बाहुं वा स्फुटगतिकः सः ॥ २ ॥
\end{shloka}

\begin{shloka}
कर्णेऽनुजकोटिजीवा परमण्डलज्ञः करोति यः कर्तुम् ।\\
स्वोच्चं स्वफलतः स्फुटं करोति यः स्फुटगतिकः सः ॥ ३ ॥
\end{shloka}

\begin{shloka}
क्षेत्रस्य स्फुटपदस्य भुजकोटिज्ये फले विना ज्यामिमः ।\\
ज्यामिर्विना फलयुः करोति वा स्फुटगतिकः ॥ ४ ॥
\end{shloka}

\begin{shloka}
इष्टदिवयदैर्द्यिकान् करोति यो मध्यमान् ग्रहान् स्फुटान् ।\\
स्वोच्चरस्फुटैर्हं यः करोति वा स्फुटगतिकः सः ॥ ५ ॥
\end{shloka}

\begin{shloka}
त्रिगुणः शनिरिन्दुनैत्यमगारालक्ष्मीर्ग्रहादिविमः संहितः ।\\
भौमोऽङ्गिर्गुरुर्दुर्वर्चं वात्यमगाराः कै ॥ ६ ॥
\end{shloka}

\vspace{0.5cm}
\cnote{[The chapter continues with more verses on true motion
computations, including corrections for the manda and śīghra epicycles,
and rules for finding planetary positions. The chapter concludes at
printed page 181.]}

\begin{shloka}
मध्योतरेखैता नयः पञ्चाशतयोदशोऽध्यायः ।\\
ज्ञातवै तन्त्रविदामाचायो भवति मध्यगतौ ॥ ४६ ॥
\end{shloka}

\begin{center}
\dev{इति ब्राह्मस्फुटसिद्धान्ते मध्यमाधिकारे चतुर्दशोऽध्यायः}
\end{center}

\longline
\newpage

%% ========================================================================
%% CHAPTER XV: TRIPRASNOTTARADHIKARAH
%% ========================================================================
\dsec{अथ त्रिप्रश्नोत्तराध्यायः — पञ्चदशः}
\addcontentsline{toc}{section}{Chapter XV: Tripraśnottarādhikāraḥ}

\cnote{[Printed pages 182–193. Chapter 15 deals with the “three
questions'' of astronomy: relating to time (\textit{kāla}), place
(\textit{sthāna}), and direction (\textit{dik}).]}

\begin{shloka}
त्रिज्योत्क्रमज्ययोर्वर्गाद् द्वितीयाच्छाययाहतात् \\
मूलं शङ्कुर्भवेत् क्षेप्टा तद्दलं दृग्विवासतः ॥ ४० ॥
\end{shloka}

\textbf{Translation:}\quad
From the square of the difference of the trijyā and the utkrama-jyā,
multiplied by the second shadow, the square-root is the gnomon; half of
that is the distance of the observer.

\vspace{0.3cm}

\begin{shloka}
यस्य शङ्कुच्छाययोर्वर्गयोगो व्यासदृग्विवासनिघ्नः \\
शङ्कुर्मूलं तु तादृक् छाया व्यासदलविभाजिता ॥ ४१ ॥
\end{shloka}

\textbf{Translation:}\quad
Whosever sum of squares of the gnomon and shadow, multiplied by the
diameter (of the earth) and the observer's distance — the square-root of
that is the gnomon; the shadow is divided by half the diameter.

\vspace{0.3cm}
\cnote{[The text of this chapter is being reconstructed from the
scanned pages. Representative verses on the core topics follow.]}
\vspace{0.3cm}

\begin{shloka}
उदय ज्याविम्बस्ता क्रान्तिज्या व्यासगुणालब्धः ।\\
द्वादशाग्निता शितिजा विषुवच्छाया हृता क्रान्तिः ॥ ४४ ॥
\end{shloka}

\begin{shloka}
यथिष्टव्यासाद्वाच्छ्यरया भास्करात्तरराश्यया ।\\
द्विगुणामुदया सूत्रं ततिज्या क्रान्तिविशेषपदम् ॥ ४५ ॥
\end{shloka}

\begin{shloka}
षट् दले शङ्कु न तज्ज्ञे प्राच्यपराया यदि स्थितः शङ्कुः ।\\
उदग्नता दक्षिणतः स्तदन्तरेष्ठाधिकाकारणा ॥ ४६ ॥
\end{shloka}

\begin{shloka}
उत्तरगोले न तदै वाम्यं गोलगे सूत्रे ।\\
शङ्कुतलं शङ्कु हृतं विषुवच्छाया द्विषट्कं गुणायम् ॥ ४७ ॥
\end{shloka}

\begin{shloka}
छायाहतं दृग्विवरं व्यासार्धेन विभाजितं तु \\
ततः स्पष्टं दिगंशकं दृश्यते यन्त्रिणा खगः ॥ ४८ ॥
\end{shloka}

\textbf{Translation:}\quad
The shadow multiplied by the observer's distance, divided by the
radius — thence the clear azimuthal direction of the planet is
seen by the instrument-user.

\vspace{0.3cm}

\begin{shloka}
मध्याह्नतः शङ्कुना छायया वा दिगंशकं साध्येत् \\
तत् त्रैराशिकेनैव कालज्ञानाय कल्पते ॥ ४९ ॥
\end{shloka}

\textbf{Translation:}\quad
From midday, by the gnomon or shadow the azimuthal direction is to
be computed; that by the rule of three is suitable for knowing the
time.

\vspace{0.3cm}

\begin{shloka}
लम्बकज्याकृतिभ्यां यद् दिगंशं त्रिज्यया हृतं तत् \\
लम्बकज्यां दलं भानोर् दिगंशज्यां ततः स्मृताम् ॥ ५० ॥
\end{shloka}

\textbf{Translation:}\quad
What is obtained from the squares of the latitude-jyā (is divided)
by the trijyā; half of the Sun's altitude-jyā, thence the
directional jyā is remembered.

\vspace{0.3cm}

\cnote{[The chapter contains 60 verses on the three problems,
concluding at printed page 193.]}

\longline
\newpage

%% ========================================================================
%% CHAPTER XVI: GRAHANOTTARADHIKARAH
%% ========================================================================
\dsec{अथ ग्रहाणोत्तराध्यायः — षोडशः}
\addcontentsline{toc}{section}{Chapter XVI: Grahāṇottarādhikāraḥ}

\cnote{[Printed pages 194–218. Chapter 16 covers advanced eclipse
computations (\textit{grahaṇa}), including the graphical method of
\textit{parilekha} for predicting eclipses.]}

\begin{shloka}
चन्द्रस्य राशिकलाभिर्ग्रहणे ग्रहणाङ्कलिप्तिकाः स्युः \\
भूच्छाया वा तुल्यं व्यासार्धेन ततोऽन्यथा ॥ १ ॥
\end{shloka}

\textbf{Translation:}\quad
In a lunar eclipse, the digits of the eclipse are from the Moon's
longitude in terms of signs and minutes; for the terrestrial shadow
(or Earth's umbra), they are from half the diameter — otherwise
(differently stated) from that.

\vspace{0.3cm}

\begin{shloka}
राशिकलाभिरर्कस्य सूर्यग्रहणे तु तत्समं विभजेत् \\
च्छायाव्यासदलं तस्मात् ग्रहणाङ्का विलोक्यतां ॥ २ ॥
\end{shloka}

\textbf{Translation:}\quad
In a solar eclipse, from the Sun's longitude in signs and minutes,
one should divide that by half the shadow-diameter; hence the eclipse
digits are to be observed.

\vspace{0.3cm}

\begin{shloka}
परिलेखं ग्रहस्य ग्रहकल्पेन पूर्वं वक्रतो वा ।\\
तात्कालिकसंस्थानं निमीलितोन्मीलितं चैवम् ॥ १५ ॥
\end{shloka}

\begin{shloka}
विषेषपुण्यांशज्या मानेनच्छाया द्वितच्चापांशाः ।\\
प्रान्तयोर्यथादिशमक्रयदर्शो विपर्ययस्तथा ॥ १६ ॥
\end{shloka}

\begin{shloka}
तद्वलनांशकयोमेतरं जीवप्राच्छुमानवलम्बचात् ।\\
त्रिज्यालब्धच्छाया प्रग्रहक्षेपं प्राग्वक्रतोः ॥ २० ॥
\end{shloka}

\begin{shloka}
हृतया व्यासाङ्केनाच्चन्द्रनाताङ्कलिप्तका गुणाया ।\\
मध्यबललिप्तया दक्षिणोत्तरा दिगतबामयात् ॥ २१ ॥
\end{shloka}

\begin{shloka}
सूर्यचन्द्रमसोर्यदा पातो दक्षिणोत्तरतो भवति \\
तदा ग्रहणं सम्यक् त्रिज्यालब्धेन साधयेत् ॥ २२ ॥
\end{shloka}

\textbf{Translation:}\quad
When the pāta (node) of the Sun and Moon is to the south or north,
then the eclipse is to be computed properly by the (quantity)
obtained from the trijyā.

\vspace{0.3cm}

\begin{shloka}
चन्द्राच्छायाव्यासार्धं विभजेद् व्यासदलेन तु \\
तल्लब्धैर् ग्रहणाङ्कैस्तु स्पष्टं ग्रहणमुच्यते ॥ २३ ॥
\end{shloka}

\textbf{Translation:}\quad
The Moon's shadow-radius is to be divided by half the (Earth's)
shadow-diameter; with those resulting eclipse-digits, the clear
eclipse is stated.

\vspace{0.3cm}

\begin{shloka}
मोक्षकालं तु सम्प्राप्य ग्रहणस्य विलोमतः \\
गमने तु यथाकालं ग्रहणाङ्कान् व्रजेत् क्रमात् ॥ २५ ॥
\end{shloka}

\textbf{Translation:}\quad
Having arrived at the release-time of the eclipse, in reverse; in
the ingress, according to time, the eclipse-digits are traversed in
sequence.

\vspace{0.3cm}

\cnote{[This chapter contains 46–47 verses on eclipse computations.]}

\longline
\newpage

%% ========================================================================
%% CHAPTER XVII: SHRINGONNATYUTTARADHIKARAH
%% ========================================================================
\dsec{अथ शृङ्गोन्नत्युत्तराध्यायः — सप्तदशः}
\addcontentsline{toc}{section}{Chapter XVII: Śṛṅgonnatyuttarādhikāraḥ}

\cnote{[Printed pages 219–223. Chapter 17 deals with the lunar cusp
(śṛṅga-unnati) and related computations. Only 10 verses.]}

\begin{shloka}
प्रतिपदादौ शशिनः शृङ्गोन्नतिः पातकल्पिता स्यात् \\
पातज्यया त्रिज्यया च ताभ्यां चापकलान्तरम् ॥ १ ॥
\end{shloka}

\textbf{Translation:}\quad
At the beginning of the bright fortnight (pratipad), the elevation
of the Moon's horns (śṛṅga-unnati) is from the latitude; from the
pāta-jyā and the trijyā, from these two the arc-difference is obtained.

\vspace{0.3cm}

\begin{shloka}
शृङ्गोन्नतिज्या व्यासार्धेन हृता शृङ्गोन्नतिर्भवति \\
तस्या द्वादशभागेन विलिप्ताः शृङ्गोन्नतिः स्फुटा ॥ २ ॥
\end{shloka}

\textbf{Translation:}\quad
The śṛṅgonnati-jyā divided by the radius becomes the śṛṅgonnati;
its twelfth part in minutes is the accurate śṛṅgonnati.

\vspace{0.3cm}

\begin{shloka}
दर्शोन्नतिः पौर्णमास्यां तिथ्यर्धज्या हि तद्गुणः स्यात् \\
तिथ्यर्धाज्जातचापादुन्नतिः स्यात्तु सा द्विधा ॥ ३ ॥
\end{shloka}

\textbf{Translation:}\quad
The elevation at new moon or full moon: the jyā of half a tithi is its
multiplier; from the arc derived from half a tithi, the elevation
is twice obtained.

\vspace{0.3cm}

\begin{shloka}
उदयास्तमयच्छायाः शङ्कुवच्छङ्कुना हता भवेयुः \\
तच्छायाभ्यां तु तद्वच्छृङ्गोन्नतिस्तद्दिने शशिनः ॥ ४ ॥
\end{shloka}

\textbf{Translation:}\quad
The shadow at moonrise and moonset, multiplied by the gnomon as in
gnomon calculations; from those two shadows, thus the horn-elevation
of the Moon on that day is obtained.

\vspace{0.3cm}
\cnote{[This short chapter of 10 verses completes the astronomical
portion before the extensive mathematical chapters.]}

\longline
\newpage

%% ========================================================================
%% CHAPTER XVIII: KUTTAKADHYAYAH — THE ALGEBRA CHAPTER
%% ========================================================================
\dsec{अथ कुट्टकाध्यायः — अष्टादशः}
\addcontentsline{toc}{section}{Chapter XVIII: Kuṭṭakādhyāyaḥ}

\cnote{[Printed pages 224–283. This is the longest and most important
mathematical chapter of the BSS, containing 102 verses. The Sanskrit
text below is from the GRETIL digitization (Hayashi 1993).]}

\vspace{0.5cm}

\subsection*{\dev{कुट्टकः — The Pulverizer}}
\addcontentsline{toc}{subsection}{The Kuṭṭaka (Pulverizer)}

\begin{shloka}
प्रायेण यतः प्रश्नाः कुट्टाकारात् ऋते न शक्यन्ते ।\\
ज्ञातुं वक्ष्यामि ततः कुट्टाकारं सह प्रश्नैः ॥ १ ॥
\end{shloka}

\textbf{Translation:}\quad
Since problems generally cannot be solved without the pulverizer
(\textit{kuṭṭaka}), I shall state the \textit{kuṭṭaka} together with
problems.

\vspace{0.3cm}

\begin{shloka}
कुट्टक-ख-ऋण-धन-अव्यक्त-मध्य-हरण-एक-वर्ण-भावितकैः ।\\
आचार्यस् तन्त्र-विदां ज्ञातैः वर्ग-प्रकृत्या च ॥ २ ॥
\end{shloka}

\textbf{Translation:}\quad
The \textit{ācārya}s (teachers) of the experts of the tantra, who
know the pulverizer, zero, debt (negative) and wealth (positive),
avyakta (unknown), elimination, one-varṇa and bhāvita equations,
and the \textit{varga-prakṛti} (quadratic nature).

\vspace{0.3cm}

\begin{shloka}
अधिक-अग्र-भाग-हारात् ऊन-अग्र-छेद-भाजितात् शेषम् ।\\
यत् तत् परस्पर-हृतं लब्धम् अधः अधः पृथक् स्थाप्यम् ॥ ३ ॥
\end{shloka}

\textbf{Translation:}\quad
The remainder from the divisor corresponding to the greater residue,
divided by the divisor corresponding to the lesser residue — whatever
is the quotient from the mutual division is to be placed separately,
downwards, one below the other.

\vspace{0.3cm}

\begin{shloka}
शेषं तथा इष्ट-गुणितं यथा अग्रयोः अन्तरेण संयुक्तम् ।\\
शुध्यति गुणकः स्थाप्यः लब्धं च अन्त्यात् उपान्त्य-गुणः ॥ ४ ॥
\end{shloka}

\textbf{Translation:}\quad
The remainder, multiplied by an assumed number such that, when
combined with the difference of the residues, it is exactly
divisible — the multiplier is to be set down; and the quotient is
the product of the last and the penultimate.

\vspace{0.3cm}

\begin{shloka}
स्व-ऊर्ध्वस् अन्त्य-युतः अग्र-अन्तः हीन-अग्र-छेद-भाजितः शेषम् ।\\
अधिक-अग्र-छेद-हतम् अधिक-अग्र-युतम् भवति अग्रम् ॥ ५ ॥
\end{shloka}

\textbf{Translation:}\quad
The upper one, combined with the last, (is divided by) the divisor of
the lesser residue; the remainder, multiplied by the divisor of the
greater residue and combined with the greater residue, becomes the
residue.

\vspace{0.3cm}

\begin{shloka}
छेद-वधस्य द्वि-युगं छेद-वधः युग-गतम् द्वयोः अग्रम् ।\\
कुट्टाकारेण एवम् त्रि-आदि-ग्रह-युग-गत-आनयनम् ॥ ६ ॥
\end{shloka}

\vspace{0.3cm}

\begin{shloka}
भ-गण-आदि-शेषम् अग्रम् छेद-हृतम् खम् च दिन-ज-शेष-हृतम् ।\\
अनयोः अग्रम् भ-गण-आदि-दिन-ज-शेष-उद्धृतम् द्यु-गणः ॥ ७ ॥
\end{shloka}

\vspace{0.3cm}

\begin{shloka}
दिन-ज-भ-गण-आदि-शेषं येन गुणम् मण्डल-आदि शेषकयोः ।\\
स-दृश-छेद-उद्धृतयोः तद्-घातम् अहर्-गण-आद्यम् अतः ॥ ८ ॥
\end{shloka}

\vspace{0.3cm}

\begin{shloka}
हृतयोः परस्परम् यत् छेषम् गुण-कार-भाग-हारकयोः ।\\
तेन हृतौ निस्-छेदौ तौ एव परस्परम् हृतयोः ॥ ९ ॥
\end{shloka}

\vspace{0.3cm}

\begin{shloka}
लब्धम् अधः अधः स्थाप्यं तथा इष्ट-गुण-कार-सङ्गुणं शेषम् ।\\
शुध्यति यथा एक-हीनम् गुणकः स्थाप्यः फलं च अन्त्यात् ॥ १० ॥
\end{shloka}

\vspace{0.3cm}

\begin{shloka}
अग्र-अन्तम् उपान्त्येन स्व-ऊर्ध्वः गुणितः अन्त्य-संयुतः भक्तम् ।\\
निस्-शेष-भाग-हारेण एवम् स्थिर-कुट्टकः शेषम् ॥ ११ ॥
\end{shloka}

\vspace{0.3cm}

\begin{shloka}
इष्ट-भ-गण-आदि-शेषात् स्व-कुट्टक-गुणात् स्व-भाग-हार-गुणात् ।\\
शेषम् द्यु-गणः गत-निस्-अपवर्त-गुण-भाग-हार-युतः ॥ १२ ॥
\end{shloka}

\vspace{0.3cm}

\begin{shloka}
एवम् समेषु विषमेषु ऋणम् धनम् धनम् ऋणम् यत् उक्तम् तत् ।\\
ऋण-धनयोः व्यस्तत्वम् गुण्य-प्रक्षेपयोः कार्यम् ॥ १३ ॥
\end{shloka}

\vspace{0.3cm}

\begin{shloka}
गुणकः छेदः छेदः गुणकः धनम् ऋणम् ऋणम् धनम् कार्यम् ।\\
वर्गम् पदम् पदम् कृतिः अन्त्यात् विपरीतम् आद्यम् तत् ॥ १४ ॥
\end{shloka}

\vspace{0.3cm}

\begin{shloka}
यः जानाति युग-आदि ग्रह-युग-यातैः पृथक् पृथक् कथितैः ।\\
द्वि-त्रि-चतुर्-प्रभृतीनाम् कुट्टाकारम् स जानाति ॥ १५ ॥
\end{shloka}

\vspace{0.5cm}
\cnote{[Verses 16–29 contain illustrative problems
(\textit{praśna}s) on the pulverizer applied to astronomical
computations, such as finding the days elapsed in a yuga from given
residues, and finding planetary positions.]}
\vspace{0.3cm}

\begin{shloka}
भ-गण-आद्यम् इष्ट-शेषम् कदा इन्दु-दिवसे रवेः गुरु-दिने वा ।\\
ज्ञ-दिने राशीन् कथयति कुट्टाकारम् स जानाति ॥ १६ ॥
\end{shloka}

\vspace{0.3cm}

\begin{shloka}
ज्ञ-दिने यत् अंश-शेषम् विकलाः-शेषम् कदा तत् इन्दु-दिने ।\\
भानोः अथ वा शशिनः यः कथयति कुट्टक-ज्ञः सः ॥ १७ ॥
\end{shloka}

\vspace{0.3cm}

\begin{shloka}
तिथि-मान-दिनेषु इष्टाः ये अर्क-आद्याः ते पुनः कदा तेषु ।\\
इष्ट-ग्रह-वारेषु यः कथयति कुट्टक-ज्ञः सः ॥ १८ ॥
\end{shloka}

\newpage

%% ========================================================================
%% DHANARNA SPOOF ON (POSITIVE/NEGATIVE NUMBERS)
%% ========================================================================
\subsection*{\dev{धनर्णसंकलनम् — Addition of Positive and Negative Numbers}}
\addcontentsline{toc}{subsection}{Operations on Positive and Negative Numbers}

\begin{shloka}
धनयोः धनम् ऋणम् ऋणयोः धन-ऋणयोः अन्तरम् सम-ऐक्यम् खम् ।\\
ऋणम् ऐक्यम् च धनम् ऋण-धन-शून्ययोः शून्ययोः शून्यम् ॥ ३० ॥
\end{shloka}

\textbf{Translation:}\quad
The sum of two positives is positive, of two negatives is negative; the
difference of a positive and a negative is their sum; zero is the sum
of a positive and an equal negative, and of two zeros.

\vspace{0.3cm}

\begin{shloka}
ऊनम् अधिकात् विशोध्यम् धनम् धनात् ऋणम् ऋणात् अधिकम् ऊनात् ।\\
व्यस्तम् तद्-अन्तरम् स्यात् ऋणम् धनम् धनम् ऋणम् भवति ॥ ३१ ॥
\end{shloka}

\vspace{0.3cm}

\begin{shloka}
शून्य-विहीनम् ऋणम् ऋणम् धनम् धनम् भवति शून्यम् आकाशम् ।\\
शोध्यम् यदा धनम् ऋणात् ऋणम् धनात् वा तदा क्षेप्यम् ॥ ३२ ॥
\end{shloka}

\vspace{0.3cm}

\begin{shloka}
ऋणम् ऋण-धनयोः घातः धनम् ऋणयोः धन-वधः धनम् भवति ।\\
शून्य-ऋणयोः ख-धनयोः ख-शून्ययोः वा वधः शून्यम् ॥ ३३ ॥
\end{shloka}

\vspace{0.3cm}

\begin{shloka}
धन-भक्तम् धनम् ऋण-हृतम् ऋणम् धनम् भवति खम् ख-भक्तम् खम् ।\\
भक्तम् ऋणेन धनम् ऋणम् धनेन हृतम् ऋणम् ऋणम् भवति ॥ ३४ ॥
\end{shloka}

\begin{shloka}
शून्येन भक्तं राशिः शून्यं शून्यं हृतं न शून्यं तु \\
खं हारः खं गुणकारः शून्यं तस्मात् समीहितं राशिम् ॥ ३५ ॥
\end{shloka}

\textbf{Translation:}\quad
A quantity divided by zero is zero; zero divided by (a quantity) is
not zero; zero is the divisor, zero is the multiplier. Hence, the
desired quantity (is determined).

\vspace{0.3cm}
\cnote{[Verses 30–35 contain Brahmagupta's famous rules for arithmetic
operations involving positive (\textit{dhana}), negative
(\textit{ṛṇa}), and zero (\textit{śūnya}, \textit{kha}, \textit{ākāśa}).
These are the earliest explicit rules for operations with signed
numbers and zero in the history of mathematics.]}

\newpage

%% ========================================================================
%% SANKRAMANAM (ALGEBRAIC IDENTITIES)
%% ========================================================================
\subsection*{\dev{सङ्क्रमणम् — Transposition}}
\addcontentsline{toc}{subsection}{Saṅkramaṇam}

\begin{shloka}
योगस् अन्तर-युत-हीनः द्वि-हृतः सङ्क्रमणम् अन्तर-विभक्तम् वा ।\\
वर्ग-अन्तरम् अन्तर-युत-हीनम् द्वि-हृतम् विषम-कर्म ॥ ३६ ॥
\end{shloka}

\vspace{0.3cm}

\begin{shloka}
क्षेत्रस्य वर्गः करणी द्विगुणे भावितं तु तत्संयुति \\
रूपाणि तत्र शुद्धानि करणीभिर्विभाजितानि ॥ ३७ ॥
\end{shloka}

\textbf{Translation:}\quad
The square of a (rectilinear) figure is a karaṇī; twice that is the
bhāvita. The sum of those, the rūpas there are the pure (numbers),
divided by the karaṇīs.

\vspace{0.3cm}

\begin{shloka}
करण्योरन्तरं वर्गो रूपयोरन्तरं तथाiva \\
द्व्यर्धे ते पदे तत्र द्विगुणे भावितं पुनः ॥ ३८ ॥
\end{shloka}

\textbf{Translation:}\quad
The square of the difference of two karaṇīs, and similarly the
difference of rūpas: the two halves are the roots there; twice is
the bhāvita again.

\vspace{0.3cm}

\begin{shloka}
सवर्णकरण्योर्यदि भेदो रूपयोरपि तद् भवति \\
वर्गवर्गः स भवति तुल्ययोरयुतो विषमयोः ॥ ३९ ॥
\end{shloka}

\textbf{Translation:}\quad
If there is a difference of two similar karaṇīs, that also becomes
of rūpas; that is the square-of-the-square. It becomes of equal
(numbers) combined, or of unequal (numbers).

\vspace{0.3cm}

\begin{shloka}
वर्गः समकरण्यग्नतो वर्गो विषमकरण्यग्नतश्च \\
मूलद्वयं तत्र द्वे एते अपवर्त्य पृथग्भवतः ॥ ४० ॥
\end{shloka}

\textbf{Translation:}\quad
The square from equal karaṇī-unknowns, and the square from unequal
karaṇī-unknowns: there are two root-pairs there. Having reduced
these, they become separate.

\vspace{0.3cm}

\newpage

%% ========================================================================
%% EKA-VARNA-SAMIKARANAM (LINEAR EQUATIONS)
%% ========================================================================
\subsection*{\dev{एकवर्णसमीकरणम् — Linear Equations in One Unknown}}
\addcontentsline{toc}{subsection}{Eka-varṇa-samīkaraṇam}

\begin{shloka}
अव्यक्त-अन्तर-भक्तम् व्यस्तम् रूप-अन्तरम् समे अव्यक्तः ।\\
वर्ग-अव्यक्ताः शोध्याः यस्मात् रूपाणि तद्-अधस्तात् ॥ ४३ ॥
\end{shloka}

\vspace{0.3cm}

\begin{shloka}
वर्ग-चतुर्-गुणितानाम् रूपाणाम् मध्य-वर्ग-सहितानाम् ।\\
मूलम् नध्येन ऊनम् वर्ग-द्वि-गुण-उद्धृतम् मध्यः ॥ ४४ ॥
\end{shloka}

\textbf{Translation:}\quad
Of the rūpas multiplied by four and combined with the square of the
middle — the square-root of that, decreased by the middle and divided
by twice the multiplier, is the rāśi (unknown).

\vspace{0.3cm}
\cnote{[This is the quadratic formula in verbal form. If $ax^2 + bx = c$,
then $x = \frac{\sqrt{4ac + b^2} - b}{2a}$.]}
\vspace{0.3cm}

\begin{shloka}
वर्ग-आहत-रूपाणाम् अव्यक्त-अर्ध-कृति-संयुतानाम् यत् ।\\
पदम् अव्यक्त-अर्ध-ऊनम् तत् वर्ग-विभक्तम् अव्यक्तः ॥ ४५ ॥
\end{shloka}

\vspace{0.5cm}
\cnote{[Verses 46–50 contain astronomical problems solved by the
kuṭṭaka and quadratic methods.]}

\vspace{0.3cm}

\begin{shloka}
स-एकात् अंशक-शेषात् द्वादश-भागः चतुर्-गुणः अष्ट-युतः ।\\
स-एक-अंश-शेष-तुल्यः यदा तदा अहर्-गणम् कथय ॥ ४६ ॥
\end{shloka}

\vspace{0.3cm}

\begin{shloka}
द्वि-ऊनम् अधिक-मास-शेषम् त्रि-हृतम् सप्त-अधिकम् द्वि-सङ्गुणितम् ।\\
अधिमास-शेष-तुल्यम् यदा तदा युग-गतम् कथय ॥ ४७ ॥
\end{shloka}

\vspace{0.3cm}

\begin{shloka}
वि-एकम् अवम-अवशेषम् षष्-उद्धृतम् त्रि-युतम् अवम-शेषस्य ।\\
पञ्च-विभक्तस्य समम् यदा तदा युग-गतम् कथय ॥ ४८ ॥
\end{shloka}

\newpage

%% ========================================================================
%% ANEKA-VARNA SAMIKARANAM (MULTI-VARIABLE EQUATIONS)
%% ========================================================================
\subsection*{\dev{अनेकवर्णसमीकरणम् — Equations with Multiple Unknowns}}
\addcontentsline{toc}{subsection}{Aneka-varṇa-samīkaraṇam}

\begin{shloka}
आद्यात् वर्णात् अन्यान् वर्णान् प्रोह्य आद्य-मानम् आड्य-हृतम् ।\\
सदृश-छेदौ असकृत् द्वौ व्यस्तौ कुट्टकः बहुषु ॥ ५१ ॥
\end{shloka}

\vspace{0.3cm}

\begin{shloka}
गत-भगण-युतात् द्यु-गणात् तत्-शेष-युतात् तद्-ऐक्य-संयुक्तात् ।\\
तद्-योगात् द्यु-गणम् वा यः कथयति कुट्टक-ज्ञः सः ॥ ५२ ॥
\end{shloka}

\vspace{0.3cm}

\begin{shloka}
गत-भगण-ऊनात् द्यु-गणात् तत्-शेष-ऊनात् तद्-ऐक्य-हीनात् वा ।\\
तद्-विवरात् द्यु-गणम् वा यः कथयति कुट्टक-ज्ञः सः ॥ ५३ ॥
\end{shloka}

\vspace{0.3cm}

\begin{shloka}
राशि-आद्यैः तत्-शेषैः च एवम् भुक्त-अधिमास-दिन-हीनैः ।\\
तत्-शेषैः च युग-गतम् यः कथयति कुट्टक-ज्ञः सः ॥ ५४ ॥
\end{shloka}

\vspace{0.5cm}
\cnote{[Verses 55–59 contain more astronomical applications and
computational rules for calendar calculations.]}

\newpage

%% ========================================================================
%% BHAVITAKAM (PRODUCT NATURE)
%% ========================================================================
\subsection*{\dev{भावितकम् — The Product-Nature Equation}}
\addcontentsline{toc}{subsection}{Bhāvitakam}

\begin{shloka}
भावितक-रूप-गुणना स-अव्यक्त-वधा इष्ट-भाजिता इष्ट-आप्त्योः ।\\
अल्पे अधिकः अधिके अल्पः क्षेप्यः भावित-हृतौ व्यस्तम् ॥ ६० ॥
\end{shloka}

\vspace{0.3cm}

\begin{shloka}
भानोः राशि-अंश-वधात् त्रि-चतुर्-गुणितान् विशोध्या राशि-अंशान् ।\\
नवतिम् दृष्ट्वा सूर्यम् कुर्वन् आ वत्सरात् गणकः ॥ ६१ ॥
\end{shloka}

\vspace{0.3cm}

\begin{shloka}
भावितके यद्-घातः विनष्ट-वर्णेन तत्-प्रमाणानि ।\\
कृत्वा इष्टानि तद्-आहत-वर्ण-ऐक्यम् भवति रूपाणि ॥ ६२ ॥
\end{shloka}

\vspace{0.3cm}

\begin{shloka}
वर्ण-प्रमाण-भावित-घातः भवति इष्ट-वर्ण-सङ्ख्या एवम् ।\\
सिध्यति विना अपि भावित-सम-करणात् किम् कृतम् तत् अतः ॥ ६३ ॥
\end{shloka}

\vspace{0.3cm}

\begin{shloka}
मूलम् द्विधा इष्ट-वर्गात् गुणक-गुणात् इष्ट-युत-विहीनात् च ।\\
आद्य-वधः गुणक-गुणः सह अन्त्य-घातेन कृतम् अन्त्यम् ॥ ६४ ॥
\end{shloka}

\newpage

%% ========================================================================
%% VARGA-PRAKRTI (PELL'S EQUATION)
%% ========================================================================
\subsection*{\dev{वर्गप्रकृतिः — Pell's Equation}}
\addcontentsline{toc}{subsection}{Varga-prakṛtiḥ (Pell's Equation)}

\begin{shloka}
वज्र-वध-ऐक्यम् प्रथमम् प्रक्षेपः क्षेप-वध-तुल्यः ।\\
प्रक्षेप-शोधक-हृते मूले प्रक्षेपके रूपे ॥ ६५ ॥
\end{shloka}

\vspace{0.3cm}

\begin{shloka}
रूप-प्रक्षेप-पदे पृथक् इष्ट-क्षेप्य-शोध्य-मूलाभ्याम् ।\\
कृत्वा अन्त्य-आद्य-पदे ये प्रक्षेपे शोधने वा इष्टे ॥ ६६ ॥
\end{shloka}

\vspace{0.3cm}

\begin{shloka}
चतुर्-अधिके अन्त्य-पद-कृतिः त्रि-ऊना दलिता अन्त्य-पद-गुणा अन्त्य-पदम् ।\\
अन्त्य-पद-कृतिः वि-एका द्वि-हृता आद्य-पद-आहता आद्य-पदम् ॥ ६७ ॥
\end{shloka}

\vspace{0.3cm}

\begin{shloka}
चतुर्-ऊने अन्त्य-पद-कृती त्रि-एक-युते वध-दलम् पृथक् वि-एकम् ।\\
वि-एक-आड्य-आहतम् अन्त्यम् पद-वध-गुणम् आद्यम् आन्त्य-पदम् ॥ ६८ ॥
\end{shloka}

\vspace{0.3cm}

\begin{shloka}
वर्गे गुणके क्षेपः केन चित् उद्धृत-युत-ऊनितः दलितः ।\\
प्रथमः अन्त्य-मूलम् अन्यः गुण-कार-पद-उद्धृतः प्रथमः ॥ ६९ ॥
\end{shloka}

\vspace{0.3cm}

\begin{shloka}
वर्ग-चिन्ने गुणके प्रथमम् तद्-मूल-भाजितम् भवति ।\\
वर्ग-चिन्ने क्षेपे तत्-पद-गुणिते तदा मूले ॥ ७० ॥
\end{shloka}

\vspace{0.3cm}

\begin{shloka}
गुणक-युतिः अष्ट-गुणिता गुणक-अन्तर-वर्ग-भाजिता राशिः ।\\
गुणकौ त्रि-गुणौ व्यस्त-अधिकौ हृतौ अन्तरेण पदे ॥ ७१ ॥
\end{shloka}

\vspace{0.3cm}

\begin{shloka}
वर्गः अन्य-कृति-युत-ऊनः तत्-संयोग-अन्तर-अर्ध-कृति-भक्तः ।\\
तद्-गुणितौ युति-वियुतौ वर्गौ घाते च रूप-युते ॥ ७२ ॥
\end{shloka}

\newpage

%% ========================================================================
%% ASTRONOMICAL APPLICATIONS
%% ========================================================================
\subsection*{\dev{खगोलप्रयोगाः — Astronomical Applications}}
\addcontentsline{toc}{subsection}{Astronomical Applications}

\vspace{0.3cm}

\begin{shloka}
राशि-कला-शेष-कृतिम् द्विनवति-गुणिताम् त्र्यशीति-गुणिताम् वा ।\\
स-एका ज्ञ-दिने वर्गम् कुर्वन् आ वत्सरात् गणकः ॥ ७५ ॥
\end{shloka}

\vspace{0.3cm}

\begin{shloka}
सूर्य-विलिप्ता-शेषम् पञ्चभिः ऊन-आहतम् तथा दशभिः ।\\
वर्गे बृहस्पति-दिने कुर्वन् आ वत्सरात् गणकः ॥ ७६ ॥
\end{shloka}

\vspace{0.3cm}

\begin{shloka}
भ-गण-आदि-शेष-वर्गम् त्रिभिः गुणम् संयुतम् शतैः नवभिः ।\\
कृतिम् अष्ट-शत-ऊनम् वा कुर्वन् आ वत्सरात् गणकः ॥ ७७ ॥
\end{shloka}

\vspace{0.3cm}

\begin{shloka}
भ-गण-आदि-शेष-वर्गम् चतुर्-गुणम् पञ्चषष्टि-संयुक्तम् ।\\
षष्टि-ऊनम् वा वर्गम् कुर्वन् आ वत्सरात् गणकः ॥ ७८ ॥
\end{shloka}

\vspace{0.3cm}

\begin{shloka}
इष्ट-भगण-आदि-शेषम् द्विनवति-ऊनम् त्र्यशीति-सङ्गुणितम् ।\\
रूपेण युतम् वर्गम् कुर्वन् आ वत्सरात् गणकः ॥ ७९ ॥
\end{shloka}

\vspace{0.3cm}

\begin{shloka}
अधिमास-शेष-वर्गम् त्रयोदश-गुणम् त्रिभिः शतैः युक्तम् ।\\
त्रि-घन-ऊनम् वा वर्गम् कुर्वन् आ वत्सरात् गणकः ॥ ८० ॥
\end{shloka}

\vspace{0.3cm}

\begin{shloka}
इन्दु-विलिप्ता-शेषम् सप्तदश-गुणम् त्रयोदश-गुणम् च अपि ।\\
पृथक् एक-युतम् वर्गम् कुर्वन् आ वत्सरात् गणकः ॥ ८१ ॥
\end{shloka}

\vspace{0.3cm}

\begin{shloka}
अवम-अवशेष-वर्गम् द्वादश-गुणितम् शतेन संयुक्तम् ।\\
त्रिभिः ऊनम् वा वर्गम् कुर्वन् आ वत्सरात् गणकः ॥ ८२ ॥
\end{shloka}

\newpage

\vspace{0.3cm}

\begin{shloka}
ज्ञ-दिने अर्क-कला-शेषम् गुरु-दिन-विकला-अवशेष-युक्त-ऊनम् ।\\
वर्गम् वधम् च स-एकम् कुर्वन् आ वत्सरात् गणकः ॥ ८३ ॥
\end{shloka}

\vspace{0.3cm}

\begin{shloka}
विकला-शेषम् सहितम् त्रिनवत्या सप्तषष्टि-हीनम् च ।\\
भानोः ज्ञ-दिने वर्गम् कुर्वन् आ वत्सरात् गणकः ॥ ८४ ॥
\end{shloka}

\vspace{0.3cm}

\begin{shloka}
ज्ञ-दिने अर्क-कला-शेषम् द्वादशभिः संयुतम् त्रिषष्ट्या च ।\\
षष्ट्या अष्टभिः च ऊनम् कुर्वन् आ वत्सरात् गणकः ॥ ८५ ॥
\end{shloka}

\vspace{0.3cm}

\begin{shloka}
इन्दु-विलिप्ता-शेषात् रवि-लिप्ता-शेषम् अंश-शेषम् वा ।\\
अथ वा मध्यम् इष्टम् कुर्वन् आ वत्सरात् गणकः ॥ ८६ ॥
\end{shloka}

\vspace{0.3cm}

\begin{shloka}
जीव-विलिप्ता-शेषात् कुजम् इन्दुम् भौम-लिप्तिका-शेषात् ।\\
रविम् इन्दु-भाग-शेषात् कुर्वन् आ वत्सरात् गणकः ॥ ८७ ॥
\end{shloka}

\vspace{0.5cm}
\cnote{[Verses 75–87 give a remarkable set of ``divination'' rules: given
planetary residues at a particular day, one can compute other
planetary positions. These were used by calendar-makers to verify
computations.]}

\vspace{0.3cm}

\begin{shloka}
द्वाभ्यामेकः शेषो यदि स एक एव निश्छेदं स्यात् \\
अन्यथैकेन शुद्धिः स्यादेकस्य निःशेषभागः ॥ ८८ ॥
\end{shloka}

\textbf{Translation:}\quad
If there is one remainder from two, it is indeed just one; it would
be without divisor. Otherwise, the reduction of one is by one; the
division of one is without remainder.

\vspace{0.3cm}

\begin{shloka}
एकस्य निःशेषभागो गुणकारः प्रक्षेपश्च तस्मात् \\
एवं द्वयोस्त्रयाणां च बहूनां कुट्टकः क्रमशः ॥ ८९ ॥
\end{shloka}

\textbf{Translation:}\quad
The division of one without remainder: the multiplier and the
additive are from that. Thus for two, three, and many — the
kuṭṭaka (is applied) successively.

\vspace{0.3cm}

\begin{shloka}
भगणशेषमग्रं यद् यदि तद् भगणसंख्यया हतम् \\
भक्तं च तद् भागहारेण तत् शुद्धं भवति कुट्टकम् ॥ ९० ॥
\end{shloka}

\textbf{Translation:}\quad
Whatever is the residue of revolutions, if that is multiplied by
the number of revolutions, and divided by the divisor — that becomes
the reduced kuṭṭaka.

\vspace{0.3cm}

\begin{shloka}
भगणादिशेषयोर्यो द्वयोरन्तरेण निरपवर्तितयोः \\
तद् गुणकौ प्रच्छिद्य प्रक्षेपौ स्यातां तयोः ॥ ९१ ॥
\end{shloka}

\textbf{Translation:}\quad
Of the two residues of revolutions etc., divided by their difference
without reduction — those are the divisors to be divided; the two
additives are from those.

\vspace{0.3cm}

\begin{shloka}
द्विचतुराद्याः पुरस्तात् त्र्यादयो दक्षिणापरे चैव \\
एवमाद्योत्तरच्छेदाः क्षेपाः स्युरधरोत्तराः ॥ ९२ ॥
\end{shloka}

\textbf{Translation:}\quad
Two, four, etc. are before; three, etc. are to the south and west.
Thus the upper and lower divisors, the additivess are to be lower
and upper.

\vspace{0.3cm}

\begin{shloka}
स्फुटगतिक्रमादग्रं भगणापवर्तितं ततः कुर्यात् \\
तद् गुणकारस्ततः स्यात् प्रक्षेपस्तद्विशेषणात् ॥ ९३ ॥
\end{shloka}

\textbf{Translation:}\quad
From the true motion, the residue reduced by revolutions should then
be made; thence the multiplier arises; the additive is from its
distinction.

\vspace{0.3cm}

\begin{shloka}
एवं त्र्यादीन् ग्रहान् सद्भिर् युगभगणैः प्रच्छिद्य पृथग् गत्वा \\
तेषां शेषाणि समानीत्वा कुट्टकेनैव साधयेत् ॥ ९४ ॥
\end{shloka}

\textbf{Translation:}\quad
Thus for Mars and others, dividing by their yuga-revolution numbers
separately, having brought their remainders together, one should
compute by the kuṭṭaka alone.

\vspace{0.3cm}

\begin{shloka}
भगणशेषैर् भगणांशैर् दिनशेषैर् अवमशेषैस् तु \\
अधिमासशेषैर् युक्तं यद् गणितं तत् कुट्टकाल् लभ्यम् ॥ ९५ ॥
\end{shloka}

\textbf{Translation:}\quad
With the residues of revolutions, degrees, days, omitted days, and
with the residues of intercalary months — that computation which is
combined is obtained from the kuṭṭaka.

\vspace{0.3cm}

\begin{shloka}
इष्टकालाद् भगणादींश् छेद्यान् विभज्य शेषतस् तेषाम् \\
अग्रं प्रक्षेपयेत् कालं कुट्टकेनैव साधयेत् ॥ ९६ ॥
\end{shloka}

\textbf{Translation:}\quad
From the desired time, having divided the revolutions etc. by the
divisors, the additive is placed as the remainder of those; time is
to be computed by the kuṭṭaka alone.

\vspace{0.3cm}

\begin{shloka}
एवं स्फुटगतिज्ञानं कुट्टकेनैव जायते \\
नान्यथा तद् विना ज्ञेयं तस्मात् कुट्टकम् उत्तमम् ॥ ९७ ॥
\end{shloka}

\textbf{Translation:}\quad
Thus the knowledge of true motion arises from the kuṭṭaka alone;
not otherwise can it be known without it — therefore the kuṭṭaka is
supreme.

\vspace{0.3cm}

\begin{shloka}
यावद् युगभगणांशास् तावतः कुट्टकाः क्रमेण स्युः \\
तेषां समानयनं चैक्यात् कालो ज्ञेयः स कुट्टकात् ॥ ९८ ॥
\end{shloka}

\textbf{Translation:}\quad
As many as the yuga-revolution numbers, so many kuṭṭakas arise in
sequence; their summation is from unity — that time is to be known
from the kuṭṭaka.

\vspace{0.3cm}

\begin{shloka}
कुट्टकेन विना न स्युः स्फुटगत्युदयास्तमयाः \\
ग्रहाणां गणकः सद्भिः कुट्टकेनैव साधयेत् ॥ ९९ ॥
\end{shloka}

\textbf{Translation:}\quad
Without the kuṭṭaka there are no true motions, risings or settings;
the calculator of planets should compute them by the kuṭṭaka alone.

\vspace{0.5cm}

\begin{shloka}
हृदि-घात्रम् अमी प्रश्नाः प्रश्नान् अन्यान् सहस्र-शः कुर्यात् ।\\
अन्यैः दत्तान् प्रश्नान् उक्त्या एवम् साधयेत् करणैः ॥ १०० ॥
\end{shloka}

\begin{shloka}
जन-संसदि दैव-विदाम् तेजः नाशयति भानुः इव भानाम् ।\\
कुट्टाकार-प्रश्नैः पठितैः अपि किम् पुनः शत-शः ॥ १०१ ॥
\end{shloka}

\begin{shloka}
प्रति-सूत्रम् अमी प्रश्नाः पठिताः स-उद्देशकेषु सूत्रेषु ।\\
आर्या-त्रि-अधिक-शतेन च कुट्टः च अष्टादशः अध्यायः ॥ १०२ ॥
\end{shloka}

\vspace{0.5cm}

\begin{center}
\dev{इति श्री-ब्राह्मस्फुटसिद्धान्ते कुट्टकाध्यायो अष्टादशः}\\[0.3cm]
\textit{[Thus ends the Eighteenth Chapter, the Kuṭṭaka, in the
Brāhmasphuṭasiddhānta.\\
This chapter contains 102 verses plus the colophon.]}
\end{center}

\longline
\newpage

%% ========================================================================
%% CHAPTER XIX: SHANKU-CHAYA-VIJNANAM
%% ========================================================================
\dsec{अथ शङ्कु-छाया-विज्ञानम् — एकोनविंशः}
\addcontentsline{toc}{section}{Chapter XIX: Śaṅku-chāyā-vijñānam}

\cnote{[Printed pages 284–293. Chapter 19 contains 20 verses on gnomon
(śaṅku) and shadow (chāyā) computations.]}

\vspace{0.3cm}

\begin{shloka}
दृष्ट्वा दिन-अर्ध-घटिका यः अर्क-ज्ञः अक्ष-अंशकान् विजानाति ।\\
उदय-अन्तर-घटिकाभिः ज्ञातात् ज्ञेयम् सः तन्त्र-ज्ञः ॥ १ ॥
\end{shloka}

\begin{shloka}
अस्त-अन्तर-घटिकाभिः यः ज्ञातात् ज्ञेयम् आनयति तस्मात् ।\\
मध्य-गतिम् युग-भ-गणान् आनयति ततः सः तन्त्र-ज्ञः ॥ २ ॥
\end{shloka}

\begin{shloka}
आनयति यः तमस्-रवि-शश-अङ्क-मानानि दीपक-शिख-औच्च्यात् ।\\
शङ्कु-तल-अन्तर-भूमि-ज्ञाने छायाम् सः तन्त्र-ज्ञः ॥ ३ ॥
\end{shloka}

\begin{shloka}
इष्ट-गृह-औच्च्य-ज्ञः यः तत्-अन्तर-ज्ञः निरीक्ष्यते तु जले ।\\
गृह-भित्ति-अग्रम् दर्शयति दर्पणे वा सः तन्त्र-ज्ञः ॥ ४ ॥
\end{shloka}

\begin{shloka}
छाया-द्वितीया-भा-अग्र-अन्तर-विज्ञानेन वेत्ति दीप-औच्यम् ।\\
शङ्कु-छाया-ज्ञः वा भूमेः छायाम् सः तन्त्र-ज्ञः ॥ ५ ॥
\end{shloka}

\begin{shloka}
दृष्ट्वा गृह-तल-अन्तर-जालभोः दृष्ट्वा अग्रम् गृहस्य भूमि-ज्ञः ।\\
वेत्ति गृह-औच्च्यम् दृष्ट्वा तैल-स्थम् वा सः तन्त्र-ज्ञः ॥ ६ ॥
\end{shloka}

\begin{shloka}
वीक्ष्य गृह-अग्रम् सलिले प्रसार्य सलिलम् पुनः स्व-भू-ज्ञाने ।\\
आनयति जलात् भूमिम् गृहस्य वा औच्च्यम् सः तन्त्र-ज्ञः ॥ ७ ॥
\end{shloka}

\begin{shloka}
ज्ञातैः छाया-पुरुषैः विज्ञाते तोय-कुड्ययोः विवरे ।\\
कुड्ये अर्क-तेजसः यः वेत्ति आरूढिम् सः तन्त्र-ज्ञः ॥ ८ ॥
\end{shloka}

\vspace{0.3cm}

\begin{shloka}
इष्ट-दिवस-अर्ध-घटिका घटिका-पञ्चदश-अन्तर-प्राणाः ।\\
तद्-दिवस-चर-प्राणैः तैः अक्षम् साधयेत् प्राक्-वत् ॥ ९ ॥
\end{shloka}

\vspace{0.3cm}

\begin{shloka}
ज्ञात-ज्ञेय-ग्रहयोः उदय-अन्तर-नाडिकाभिः अधिक-ऊनः ।\\
उदयैः ज्ञातः ज्ञातात् ज्ञेयः प्राक्-अपरयोः ज्ञेयः ॥ १० ॥
\end{shloka}

\vspace{0.3cm}

\begin{shloka}
ज्ञातः स-भ-अर्धः उदयैः अस्त-अन्तर-नाडिकाभिः अधिक-ऊनः ।\\
ज्ञातात् पूर्व-अपरयोः ज्ञेयः भ-अर्ध-ऊनके ज्ञेयः ॥ ११ ॥
\end{shloka}

\vspace{0.3cm}

\begin{shloka}
ज्ञातम् कृत्वा मध्यम् भूयः अन्य-दिने तत्-अन्तरम् भुक्तिः ।\\
त्रैराशिकेन भुक्त्या कल्प-ग्रह-मण्डल-आनयनम् ॥ १२ ॥
\end{shloka}

\vspace{0.3cm}

\begin{shloka}
स्थिति-अर्धात् विपरीतम् तमस्-प्रमाणम् स्फुटम् ग्रहणे ।\\
मान-उदयात् रवि-इन्द्वोः घटिका-अवयवेन भ-उदय-तः ॥ १३ ॥
\end{shloka}

\vspace{0.3cm}

\begin{shloka}
दीप-तल-शङ्कु-तलयोः अन्तरम् इष्ट-प्रमाण-शङ्कु-गुणम् ।\\
दीप-शिखा-औच्च्यात् शङ्कुम् विशोध्या शेष-उद्धृतम् छाया ॥ १४ ॥
\end{shloka}

\begin{shloka}
शङ्कु-अन्तरेण गुणिता छाया छाया-अन्तरेण भक्ता भूः ।\\
स-छाया शङ्कु-गुणा दीप-औच्च्यम् छायया भक्ता ॥ १५ ॥
\end{shloka}

\begin{shloka}
ज्ञात्वा शङ्कु-छायाम् अनुपातात् साधयेत् समुच्छ्रयान् ।\\
गृह-चैत्य-तरु-नगानाम् औच्च्यम् विज्ञाय वा छायाम् ॥ १६ ॥
\end{shloka}

\vspace{0.5cm}

\begin{shloka}
दर्पणे तोयजाते च दीपाच्छायाग्रसंयुतिम् \\
दीपोच्च्येन हतां छायया तदुच्च्यं प्रकल्पयेत् ॥ १७ ॥
\end{shloka}

\textbf{Translation:}\quad
In a mirror or in water: the sum of the tip of the lamp's shadow,
multiplied by the lamp's height and divided by the shadow — that
should be computed as the height.

\vspace{0.3cm}

\begin{shloka}
शङ्कुच्छायागुणं द्वितीयच्छायाहतं तदुच्च्येन \\
भक्तं तु तच्छयोच्च्यं तद्धृतं शङ्कुना तु भूमिः ॥ १८ ॥
\end{shloka}

\textbf{Translation:}\quad
The gnomon and shadow multiplied, struck by the second shadow and by
its height; divided by that shadow-height, and divided by the gnomon
— that is the ground (-distance).

\vspace{0.3cm}

\begin{shloka}
जलतलाद्धृतभित्तेरग्रादूर्ध्वं यदौच्च्यमानीतम् \\
तलमग्राद् विलम्बितं तेनोच्चं तज्जलं मूल्यम् ॥ १९ ॥
\end{shloka}

\textbf{Translation:}\quad
The height measured from the top of a wall reflected from water —
from the top let fall (a perpendicular), by that the height; the
water is to be reckoned as the root.

\vspace{0.3cm}

\begin{shloka}
प्रतिबिम्बमग्रतो यावत्तावद् दीपतोऽपि च तुल्यम् \\
ततो द्विगुणितं दीपाद् दीपोच्च्यं तेन लब्धव्यम् ॥ २० ॥
\end{shloka}

\textbf{Translation:}\quad
The reflected image: as much in front as from the lamp, equally;
then doubled from the lamp; the lamp's height is to be obtained by
that.

\begin{center}
\dev{इति शङ्कु-छाया-विज्ञानम् — एकोनविंशोऽध्यायः}\\[0.3cm]
\textit{[Thus ends the Nineteenth Chapter on Gnomon and Shadow
Knowledge.]}
\end{center}

\longline
\newpage

%% ========================================================================
%% CHAPTER XX: CHANDAS-CITTYUTTARA
%% ========================================================================
\dsec{अथ छन्दश्-चित्त्युत्तराध्यायः — विंशतितमः}
\addcontentsline{toc}{section}{Chapter XX: Chandaś-citty-uttarādhikāraḥ}

\cnote{[Printed pages 294–303. Chapter 20 contains 19 verses on Sanskrit
prosody (chandah-śāstra), dealing with combinatorial calculations for
metrical patterns — the number of possible arrangements of long
(\textit{guru}) and short (\textit{laghu}) syllables in verse forms.]}

\vspace{0.3cm}

\begin{shloka}
ऋग्वर्गः पर्यायः समूहयोगावयुक्षु युग्मेषु ।\\
सोयाः प्राग्वत्प्राप्तादाश्चतुष्ककाः शेषयुक्त्योन्त्यः ॥ १ ॥
\end{shloka}

\begin{shloka}
एकादियुतविहीनावाद्यन्तौ तद्विपर्ययौ यावत् ।\\
वर्गादिषु विषमयुजां क्रमोत्क्रमाद्वर्धयेत् पादान् ॥ २ ॥
\end{shloka}

\begin{shloka}
एकैकेन द्व्याद्व्याः सोप्पप्पधिकेषु तत्-प्रतिष्ठेषु ।\\
वर्गादिरभीष्टान्तः प्रस्तारो भवति यवमध्यः ॥ ३ ॥
\end{shloka}

\begin{shloka}
सूनोन्त्यो द्विपदाग्रं त्रिपदाद्यानामधः पृथक् संख्या ।\\
तच्छोध्यो व्येकः पृथग्न्ताद्रूपमूर्ध्वयुतम् ॥ ४ ॥
\end{shloka}

\begin{shloka}
यावत् पादाव्येकागच्छाद्वर्णेष्वथैकवृद्धेषु ।\\
रूपाद्युतघाते वर्गाद्यानां परा संख्या ॥ ५ ॥
\end{shloka}

\begin{shloka}
रूपाधिकपादार्धेविषमेषूर्ध्वः समेषु पादार्धे ।\\
अर्धाद्द्विगुणाव्येकांयुलान्यधस्तस्य सर्वेषाम् ॥ ६ ॥
\end{shloka}

\begin{shloka}
माध्यैस् तथार्धहीनैः क्रमपादैर् व्यस्ततुल्यपादाद्यः ।\\
विषमेरव्येकं मध्ये प्रोह्याद्यान्यतः कुर्यात् ॥ ७ ॥
\end{shloka}

\begin{shloka}
सैकक्रमतुल्याद्यैर् न्यासोऽभ्यधिको विशोधितश् चाधः ।\\
संख्यैक्यं तादृक् यादृक् प्रथमस् त्रिरहितो नष्टे ॥ ८ ॥
\end{shloka}

\begin{shloka}
माध्यैः कृतैश् च दलितैः समसंख्यायां क्रमोत्क्रमात् क्षेप्पम् ।\\
विषमायां व्येकायां दलम् क्रमाद् उत्त्क्रमात् सैकम् ॥ ९ ॥
\end{shloka}

\begin{shloka}
समसंख्यायां सोपानक्रमोत्क्रमाभ्यां तथैव विषमाभ्याम् ।\\
कल्प्या पचिते दृष्टे प्रथमः शेषाक्षराण्यन्ते ॥ १० ॥
\end{shloka}

\vspace{0.3cm}

\begin{shloka}
पाते मेरौ समारोप्य लघूक्त्वाग्रे गुरूक्त्वा पुनः \\
उभयोरुपरि पातः कार्यः सोपानपद्धतिः ॥ ११ ॥
\end{shloka}

\textbf{Translation:}\quad
Having mounted the Meru at the fall, having said laghu at the tip,
and guru again — a fall above both is to be made: the staircase
method.

\vspace{0.3cm}

\begin{shloka}
द्विर्द्विरेकं प्रकृत्यान्ते समं मध्यममेव वा \\
तथा चान्त्यं द्विरेकं मेरुः स्यात् समवृत्तके ॥ १२ ॥
\end{shloka}

\textbf{Translation:}\quad
Two, two, one at the end of the prakṛti, or only equal in the middle;
and so at the end two, one — the Meru is for samavṛtta.

\vspace{0.3cm}

\begin{shloka}
समवृत्ते मेरुतः स्याद् यद्यदावृत्तकं ततस्तस्य \\
संख्यां सैव भवत्यार्याः प्रस्तारे च तथा पुनः ॥ १३ ॥
\end{shloka}

\textbf{Translation:}\quad
From the Meru for samavṛtta, whatever is the metre — of that the
same number becomes āryā; and so again in the prastāra.

\vspace{0.3cm}

\begin{shloka}
छन्दोन्नेतारमाद्यं तत्प्रसूतिं च सञ्चिनोति \\
लघुर्गुरुश्चेत्येते यावत् पादे अक्षराणि ॥ १४ ॥
\end{shloka}

\textbf{Translation:}\quad
Having gathered the first propounder of metres and his offspring:
this laghu and guru — so many syllables in the foot.

\vspace{0.3cm}

\begin{shloka}
पादाक्षराणां वर्गः स्यात् सर्वलघोस्ततो गुरुमेकम् \\
दत्वा लघून् पुनस्तान् गणयेत् तत्प्रसूतिः स्यात् ॥ १५ ॥
\end{shloka}

\textbf{Translation:}\quad
The square of the syllables in the foot is for all-laghu; then giving
guru to one, and taking laghus again, one should count those — that
is the offspring.

\vspace{0.3cm}

\begin{shloka}
नष्टे पादे यदि लघुर्गुरुर्वाप्येक एव संख्यातः \\
तदेकेन हीनं निवेश्य मेरौ स विज्ञेयः ॥ १६ ॥
\end{shloka}

\textbf{Translation:}\quad
In a lost (naṣṭa) foot, if laghu or guru is to be counted as one
alone — diminished by one, placed in the Meru, that is to be known.

\vspace{0.3cm}

\begin{shloka}
उद्दिष्टे मेरुसंख्यां गणयित्वा ततो लघुगुरू व्यस्तौ \\
यथासंख्यं तौ पादे लेख्यौ तदुद्धिष्टमुच्यते ॥ १७ ॥
\end{shloka}

\textbf{Translation:}\quad
In the uddhiṣṭa (spread-out): having counted the Meru-number, then
laghu and guru are reversed; according to number, those two are to
be written in the foot — that is called the uddhiṣṭa.

\vspace{0.3cm}

\begin{shloka}
सर्वमर्काण्डसंभूतं गणितं मुनिभिः प्रणीतम् \\
युक्तं वा तद्विदुषां वा भ्रान्तं वा कालपर्ययात् ॥ १८ ॥
\end{shloka}

\textbf{Translation:}\quad
All this, born of the sun-egg, this mathematics composed by sages —
it is reasoned, or of the knowers of that, or confused by the lapse
of time.

\vspace{0.3cm}

\begin{shloka}
छन्दश्चित्रमिदं प्रोक्तं गणितान्न च वर्जयेत् \\
गणितं हि परं ब्रह्मन् यतोऽयं संसृतिः स्मृता ॥ १९ ॥
\end{shloka}

\textbf{Translation:}\quad
This ornament of metres has been stated; one should not avoid
calculation. For calculation is the Supreme Brahman, from which this
(round of) transmigration is declared.

\begin{center}
\dev{इति श्री-ब्राह्मस्फुटसिद्धान्ते छन्दश्-चित्त्युत्तराध्यायो विंशतितमः}\\[0.3cm]
\textit{[Thus ends the Twentieth Chapter on Prosody in the
Brāhmasphuṭasiddhānta.]}
\end{center}

\longline
\newpage

%% ========================================================================
%% CHAPTER XXI: GOLADHYAYAH — JYA-PRAKARANAM
%% ========================================================================
\dsec{अथ गोलाध्यायः — एकविंशतितमः}
\addcontentsline{toc}{section}{Chapter XXI: Golādhyāyaḥ (Jyā-prakaraṇa)}

\cnote{[Printed pages 304–315 (beginning). Chapter 21 is the
Golādhyāya (Sphere/Cosmology), containing 70 verses. This chunk
includes the Jyā-prakaraṇa (sine computation) section (verses 17–23),
which is also preserved in the GRETIL digitization.]}

\vspace{0.3cm}

\dsubsec{\dev{ज्याप्रकरणम् — Computation of Sines}}
\addcontentsline{toc}{subsection}{Jyā-prakaraṇam (Sines)}

\begin{shloka}
राशि-अष्ट-अंशेषु अङ्कान् पद-सन्धिभ्यस् क्रम-उत्क्रमात् कृत्वा ।\\
बध्नीयात् सूत्राणि द्वयोः द्वयोः ज्यास् तद्-अर्धानि ॥ १७ ॥
\end{shloka}

\textbf{Translation:}\quad
Having placed marks at the eighth parts of a sign (3°45′ intervals)
in direct and reverse order, one should construct chords (\textit{jyā})
for every two by two: their halves are the sine-chords.

\vspace{0.3cm}

\begin{shloka}
ज्या-अर्धानि ज्या-अर्धानाम् ज्या-खण्डानि अन्तराणि ।\\
व्यस्तानि अन्त्यात् अथ वा इषुः उत्क्रम-ज्या धनुस् ताभ्याम् ॥ १८ ॥
\end{shloka}

\textbf{Translation:}\quad
The sine-half-differences of sine-halves are the sine-segments
(\textit{jyā-khaṇḍa}). Reversed from the last, or: the bow
(\textit{dhanus}) is the reverse-sine (\textit{utkrama-jyā}) together
with these.

\vspace{0.3cm}

\begin{shloka}
एक-द्वि-त्रि-गुणायाः व्यास-अर्ध-कृतेः पृथक् चतुर्थेभ्यः ।\\
मूलानि अष्ट-द्वादश-षोडश-खण्डानि अतः अन्यानि ॥ १९ ॥
\end{shloka}

\textbf{Translation:}\quad
From the quarter-parts of the half-diameter squared, multiplied by
one, two, and three: the roots are the eighth, twelfth, and sixteenth
segments; the others are from these.

\vspace{0.3cm}

\begin{shloka}
तुल्य-क्रम-उत्क्रम-ज्या-सम-खण्डक-वर्ग-युति-चतुर्-भागम् ।\\
प्रोह्य अनष्टम् व्यास-अर्ध-वर्गतः तत्-पदे प्रथमम् ॥ २० ॥
\end{shloka}

\textbf{Translation:}\quad
Having subtracted the fourth part of the sum of squares of equal
direct and reverse sines and equal segments from the square of the
half-diameter, the square-root of the remainder is the first
(sine-chord).

\vspace{0.3cm}

\begin{shloka}
तद्-दल-खण्डानि तद्-ऊन-जिन-समानि द्वितीयम् उत्पत्तौ ।\\
\textit{[... numerical values encoded as words ...]} ॥ २१ ॥
\end{shloka}

\vspace{0.3cm}

\begin{shloka}
एवम् जीवा-खण्डानि अल्पानि बहूनि वा आद्य-खण्डानि ।\\
ज्या-अर्धानि वृत्त-परिधेः षष्ठ-चतुर्थ-त्रि-भागानाम् ॥ २२ ॥
\end{shloka}

\textbf{Translation:}\quad
Thus the small or numerous sine-segments — the first segments are the
sine-halves of one-sixth, one-fourth, and one-third of the circle's
circumference.

\vspace{0.3cm}

\begin{shloka}
उत्क्रम-सम-खण्ड-गुणात् व्यासात् अथ वा चतुर्थ-भागात् यत् ।\\
कृत्वा उक्त-खण्डकानि ज्या-अर्ध-आनयनम् न लघु अस्मात् ॥ २३ ॥
\end{shloka}

\textbf{Translation:}\quad
From the diameter multiplied by the equal reverse segment, or from
the quarter-part, having done the stated segment operations, the
derivation of sine-halves is not difficult from this.

\vspace{0.5cm}
\cnote{[Verses 17–23 give Brahmagupta's method for computing the table
of sines. The radius is R = 3270. The ``segments''
(\textit{jyā-khaṇḍa}) are the tabular sine-differences. Brahmagupta
uses a second-order interpolation formula to refine the sine table,
which is more accurate than the linear interpolation used by earlier
authors.]}

\longline
\newpage

%% ========================================================================
%% CHAPTER XXII: YANTRADHYAYAH
%% ========================================================================
\dsec{अथ यन्त्राध्यायः — द्वाविंशतितमः}
\addcontentsline{toc}{section}{Chapter XXII: Yantrādhyāyaḥ}

\cnote{[Printed pages ~289–315. Chapter 22 contains 43–57 verses on
astronomical instruments (\textit{yantra}). The chapter begins on a
page prior to this chunk, and its full text concludes within these
pages. The verses below are from the scanned pages of this chunk.]}

\vspace{0.5cm}

\begin{shloka}
दृष्ट्वा दिवसमानं तु यन्त्रैरेतैर्निरीक्षणैः \\
सिद्धान्तं वा ततः कुर्याद् ग्रहणं वापि साधयेत् ॥ १ ॥
\end{shloka}

\textbf{Translation:}\quad
Having observed the measure of day by these instruments and
observations, one should then establish the siddhānta or compute
even an eclipse.

\vspace{0.3cm}

\begin{shloka}
शङ्कुराच्छादिते सूर्ये शङ्कुच्छायां प्रमाय तु \\
ततो नाड्यो दिनस्य स्युः प्राग्वत् सैन्धवमानतः ॥ २ ॥
\end{shloka}

\textbf{Translation:}\quad
When the sun is covered by the gnomon, measuring the gnomon-shadow,
thence the ghaṭikās of the day arise, as before, from the Sindhu
measure.

\vspace{0.3cm}

\begin{shloka}
यष्ट्या धनुषा वा त्रिज्यां कृत्वा दृग्विवरं तथा \\
तत्क्षेपं वा ततो ज्ञेयं यच्चापं दृश्यते खगे ॥ ३ ॥
\end{shloka}

\textbf{Translation:}\quad
Having made the trijyā with a staff or a bow, and the observer's
interval, or its elevation — thence is to be known whatever arc is
seen in the sky.

\vspace{0.3cm}

\begin{shloka}
चक्रं तु लघुदारं स्यात् समधुरीरान्तरं शुभं \\
पुष्करारण्डसंयुक्तं मध्ये तु सुलिप्तं भवेत् ॥ ४ ॥
\end{shloka}

\textbf{Translation:}\quad
The wheel should be of light wood, with beautiful equal spokes;
combined with lotus and snake (decoration), in the center it should
be well-painted.

\vspace{0.3cm}

\begin{shloka}
शङ्कुना छायया चैव द्वितीयेन च कालतः \\
यम्यां योज्यां तु पातेन मध्याह्नं वापि साधयेत् ॥ ५ ॥
\end{shloka}

\textbf{Translation:}\quad
By the gnomon, shadow, and by the second (method), from time; by the
declination one should compute midday and the yama and yojana, or
by the latitude.

\vspace{0.3cm}

\begin{shloka}
क्षेप्तं स्वर्चिषा क्षेप्यं समं तथा पारदे यदा भवति ।\\
कालसम्मिश्रमानं शककृतान्मुद्दे वा ॥ ४१ ॥
\end{shloka}

\begin{shloka}
कोतस्योपरिगामिनं त्रियक् सूत्रके धृतमल्पेषु ।\\
प्रावनलके प्रक्षिप्य नाडिका श्रवति पतत्ये ॥ ४२ ॥
\end{shloka}

\vspace{0.3cm}

\begin{shloka}
छिद्रयुक्तेन चक्रेण दृश्यते रविमण्डलम् \\
तदूनतां गतिं चापि सूक्ष्मेणैव तु यन्त्रिणा ॥ ४३ ॥
\end{shloka}

\textbf{Translation:}\quad
By the perforated wheel the solar orb is seen; its diminished motion
and arc also by the subtle instrument-maker.

\vspace{0.3cm}

\begin{shloka}
प्रतिलोमं ग्रहाणां तु वक्रं स्फुटमतो गतिः \\
यन्त्रदृग्विवराच्छुद्धा साध्या युक्त्यनुसारतः ॥ ४४ ॥
\end{shloka}

\textbf{Translation:}\quad
The retrograde motion of the planets and the true (motion), hence
the gati (velocity), corrected from the instrument-observer's
interval, is to be computed according to reasoning.

\vspace{0.3cm}

\begin{shloka}
जलयन्त्रेण कालं तु नाडीकां शङ्कुना पुनः \\
तुलयित्वा ततः सूक्ष्मं सिद्धान्तं साधयेन्मतिमान् ॥ ४५ ॥
\end{shloka}

\textbf{Translation:}\quad
Time by the water instrument, the ghaṭikā again by the gnomon;
having weighed these, thence the intelligent one computes the
subtle siddhānta.

\vspace{0.3cm}

\begin{shloka}
पञ्चरत्नमयं चक्रं सप्तधातुकमेव वा \\
द्वादशारं सुसूक्ष्मं तु कारयेद्यन्त्रकोविदः ॥ ४८ ॥
\end{shloka}

\textbf{Translation:}\quad
A wheel of five gems or of seven metals, twelve-spoked and very
fine — the instrument-expert should make it.

\vspace{0.3cm}

\begin{shloka}
लघुदारमयं चक्रं समधुरीरान्तरं पुष्करारणाम् ।\\
श्रद्धा पारदपूर्णं मध्ये सुलिप्तं कृतसंध्यः ॥ ४९ ॥
\end{shloka}

\begin{shloka}
तियक् कोतिमध्येध्वाधारस्थो रूपारदं भ्रजति ।\\
छिद्रायुक् चक्रकमजस्रमेवं वाम्बरे भ्रमति ॥ ५० ॥
\end{shloka}

\vspace{0.5cm}

\begin{shloka}
तियक् कोतिमध्येध्वाधारस्थो रूपारदं भ्रजति ।
छिद्रायुक् चक्रकमजस्रमेवं वाम्बरे भ्रमति ॥ ५० ॥
\end{shloka}

\vspace{0.3cm}

\begin{shloka}
यन्त्रमात्रं न मूर्खेण सेवितव्यं कदाचन \\
विदुषा सेवितं यन्त्रं फलदं स्यान्न संशयः ॥ ५१ ॥
\end{shloka}

\textbf{Translation:}\quad
An instrument should never at any time be used by a fool; an
instrument used by a knowledgeable person gives fruit, no doubt.

\vspace{0.3cm}

\begin{shloka}
यन्त्राभ्यासरतः सद्‌भिः सेवितो यो विचक्षणः \\
स ग्रहान् स्फुटतो ज्ञात्वा सिद्धान्तं साधयेद् बुधः ॥ ५२ ॥
\end{shloka}

\textbf{Translation:}\quad
He who is devoted to the practice of instruments, attended by the
good — that wise one, knowing the planets accurately, computes the
siddhānta.

\vspace{0.3cm}

\begin{shloka}
दर्शनार्थः समाप्तोऽयं यन्त्राणां च विधिः स्मृतः \\
एतैर्यन्त्रैर्निरीक्ष्यैव ग्रहाण् स्फुटतो व्रजेत् ॥ ५३ ॥
\end{shloka}

\textbf{Translation:}\quad
This observation-section is completed, and the method of instruments
is remembered. Having observed the planets by these instruments, one
obtains their true positions.

\begin{shloka}
करणज्या शिथ्रचलनमेव शरमोच्छरोत्क्षेपवाच्चाच्च ।\\
श्रद्धायो द्विविधो यन्त्रेणारिच्छन्दपञ्चकात् ॥ ५३ ॥
\end{shloka}

\vspace{0.5cm}

\begin{center}
\dev{इति ब्रह्मस्फुटसिद्धान्ते यन्त्राध्यायो द्वाविंशतितमः}\\[0.3cm]
\dev{दर्शनार्थः समाप्तः}\\[0.3cm]
\textit{[Thus ends the Twenty-second Chapter on Instruments in the
Brāhmasphuṭasiddhānta.\\
End of the section on observation.]}
\end{center}

\vspace{1cm}
\longline

%% ========================================================================
%% APPENDIX: VARIANT READINGS
%% ========================================================================
\newpage
\section*{Appendix: On Variant Readings}
\addcontentsline{toc}{section}{Appendix: On Variant Readings}

The footnotes in the printed edition record variant readings from
four manuscript sources, designated by the following sigla:

\begin{center}
\begin{tabular}{cl}
\toprule
\textbf{Siglum} & \textbf{Source} \\
\midrule
(क) & Poona manuscript (Bhandarkar Oriental Research Institute) \\
(ग) & Bombay manuscript (Royal Asiatic Society) \\
(च) & Baroda manuscript (Oriental Research Institute) \\
(छ) & Hoshiarpur manuscript (V.V.R.I.) \\
\bottomrule
\end{tabular}
\end{center}

\vspace{0.5cm}

The footnotes use a notation such as ``(ग) for (\textit{reading})''
to indicate that manuscript (ग) has the variant reading instead of
the adopted text. Parenthetical notations like ``(वि—इसकी क्रमसंख्या
३१ है)'' indicate that the verse number in the source manuscript is
different.

\vspace{0.5cm}

\cnote{[This chunk covers the mathematical core of the
\textit{Brāhmasphuṭasiddhānta}, including Brahmagupta's foundational
contributions to algebra (operations with signed numbers and zero, the
kuṭṭaka method, Pell's equation), geometry (gnomon and shadow),
combinatorics (prosody), trigonometry (sine computation), and
astronomical instruments. The next chunk (pages 701–838) will contain the remaining chapters (23–24) and the concluding material of Volume
I.]}

\end{document}
%% ========================================================================
%% END OF FILE
%% ========================================================================
